\SourceLeaf{185}
\Needspace{8\baselineskip}
\Article{193}
\textsc{Satz.} \textit{Sind die reduzierten Formen $f$, $F$ eigentlich äquivalent, so ist die eine in der Periode der anderen enthalten.}

Sei $f=(a,b,-a')$, $F=(A,B,-A')$, die Determinante dieser Formen sei $D$, und die erste gehe durch die eigentliche Substitution $\mathfrak{A},\mathfrak{B},\mathfrak{C},\mathfrak{D}$ in die zweite über. Dann behaupte ich: Wenn wir die Periode der Form $f$ suchen und wie in Artikel 188 die nach beiden Seiten unendliche Folge der reduzierten Formen und der Transformationen von $f$ in diese bestimmen, so ist \textit{entweder} $+\mathfrak{A}$ einem Glied der Folge $\mathinner{\ldotp\ldotp\ldotp}{}''\alpha,{}'\alpha,\alpha,\alpha',\alpha''\mathinner{\ldotp\ldotp\ldotp}$ gleich, und wenn wir dieses $=\alpha^m$ setzen, ist $+\mathfrak{B}$ $=\beta^m$, $+\mathfrak{C}=\gamma^m$, $+\mathfrak{D}=\delta^m$; \textit{oder} $-\mathfrak{A}$ ist einem Glied $\alpha^m$ gleich, und $-\mathfrak{B},-\mathfrak{C},-\mathfrak{D}$ sind der Reihe nach $=\beta^m,\gamma^m,\delta^m$ (wobei $m$ auch einen negativen Index bezeichnen kann). In beiden Fällen ist $F$ offenbar mit $f^m$ identisch.

\textit{Beweis.} I. Wir haben die vier Gleichungen
\begin{align*}
a\mathfrak{A}\mathfrak{A}+2b\mathfrak{A}\mathfrak{C}-a'\mathfrak{C}\mathfrak{C}&=A \quad \mathinner{\ldotp\ldotp\ldotp\ldotp\ldotp\ldotp\ldotp} \tag*{[1]}\\
a\mathfrak{A}\mathfrak{B}+b(\mathfrak{A}\mathfrak{D}+\mathfrak{B}\mathfrak{C})-a'\mathfrak{C}\mathfrak{D}&=B \quad \mathinner{\ldotp\ldotp\ldotp\ldotp\ldotp\ldotp} \tag*{[2]}\\
a\mathfrak{B}\mathfrak{B}+2b\mathfrak{B}\mathfrak{D}-a'\mathfrak{D}\mathfrak{D}&=-A' \quad \mathinner{\ldotp\ldotp\ldotp\ldotp\ldotp\ldotp} \tag*{[3]}\\
\mathfrak{A}\mathfrak{D}-\mathfrak{B}\mathfrak{C}&=1 \quad \mathinner{\ldotp\ldotp\ldotp\ldotp\ldotp\ldotp\ldotp\ldotp\ldotp\ldotp} \tag*{[4]}
\end{align*}
Wir betrachten \textit{zuerst} den Fall, dass eine der Zahlen $\mathfrak{A}\ \mathfrak{B},\mathfrak{C},\mathfrak{D}=0$ ist.

$1^\circ$ Ist $\mathfrak{A}=0$, so folgt aus [4] $\mathfrak{B}\mathfrak{C}=-1$, also $\mathfrak{B}=\pm1$, $\mathfrak{C}=\mp1$. Dann ergibt [1] $-a'=A$ und [2] $-b\pm a'\mathfrak{D}=B$ oder $B\equiv-b\ (\mathrm{mod.}\ a'\text{ oder }A)$. Daraus folgt, dass die Form $(A,B,-A')$ der Form $(a,b,-a')$ nach rechts benachbart ist. Da die erste reduziert ist, muss sie mit $f'$ identisch sein. Also ist $B=b'$, und aus [2] folgt $b+b'=-a'\mathfrak{C}\mathfrak{D}=\pm a'\mathfrak{D}$; wegen $\frac{b+b'}{-a'}=h'$ erhalten wir $\mathfrak{D}=\mp h'$. Wir schließen, dass $\mp\mathfrak{A},\mp\mathfrak{B},\mp\mathfrak{C},\mp\mathfrak{D}$ der Reihe nach $=0,-1,+1,h'$ oder $=\alpha',\beta',\gamma',\delta'$ sind.

$2^\circ$ Ist $\mathfrak{B}=0$, so ergibt [4] $\mathfrak{A}=\pm1$, $\mathfrak{D}=\pm1$; [3] ergibt $a'=A'$; und [2] ergibt $b\mp a'\mathfrak{C}=B$ oder $b\equiv B\ (\mathrm{mod.}\ a')$. Da sowohl $f$ als auch $F$ reduzierte Formen sind, liegen sowohl $b$ als auch $B$ zwischen $\surd D$ und $\surd D\mp a'$ (je nachdem, ob $a'$ positiv oder negativ ist, Artikel 185, 5). Daher muss $b=B$ und $\mathfrak{C}=0$ sein. Somit sind die Formen $f$, $F$ identisch, und es gilt der Reihe nach $\pm\mathfrak{A},\pm\mathfrak{B},\pm\mathfrak{C},\pm\mathfrak{D}=1,0,0,1=\alpha,\beta,\gamma\ \delta$.

$3^\circ$ Ist $\mathfrak{C}=0$, so ergibt [4] $\mathfrak{A}=\pm1$, $\mathfrak{D}=\pm1$; [1] ergibt $a=A$; und [2] ergibt $\pm a\mathfrak{B}+b=B$ oder $b\equiv B\ (\mathrm{mod.}\ a)$. Da aber sowohl $b$ als auch $B$ zwischen $\surd D$
\SourceLeaf{186}
und $\surd D\mp a$ liegen, muss $B=b$ und $\mathfrak{B}=0$ sein. Dieser Fall unterscheidet sich also nicht vom vorhergehenden.

$4^\circ$ Ist $\mathfrak{D}=0$, so ergibt [4] $\mathfrak{B}=\pm1$, $\mathfrak{C}=\mp1$; [3] ergibt $a=-A'$; und [2] ergibt $\pm a\mathfrak{A}-b=B$ oder $B\equiv-b\ (\mathrm{mod.}\ a)$. Daher ist $F$ der Form $f$ nach links benachbart und somit mit der Form ${}'f$ identisch. Wegen $\frac{{}'b+b}{a}=h$ und $B={}'b$ erhalten wir $\pm\mathfrak{A}=h$. Wir schließen, dass $\pm\mathfrak{A},\pm\mathfrak{B},\pm\mathfrak{C},\pm\mathfrak{D}$ der Reihe nach $=h,1,-1,0$, $={}'\alpha,{}'\beta,{}'\gamma,{}'\delta$ sind.

Es bleibt der Fall, dass keine der Zahlen $\mathfrak{A},\mathfrak{B},\mathfrak{C},\mathfrak{D}=0$ ist. Nach dem Hilfssatz in Artikel 190 haben dann die Größen $\frac{\mathfrak{A}}{\mathfrak{C}}$, $\frac{\mathfrak{B}}{\mathfrak{D}}$, $\frac{\mathfrak{C}}{\mathfrak{A}}$, $\frac{\mathfrak{D}}{\mathfrak{B}}$ dasselbe Vorzeichen. Daraus ergeben sich zwei Fälle, denn dieses Vorzeichen kann entweder mit dem von $a$, $a'$ übereinstimmen oder ihm entgegengesetzt sein.

II. Haben $\frac{\mathfrak{A}}{\mathfrak{C}}$, $\frac{\mathfrak{B}}{\mathfrak{D}}$ dasselbe Vorzeichen wie $a$, so liegt die Größe $\frac{\surd D-b}{a}$ (die wir mit $L$ bezeichnen werden) zwischen diesen Brüchen (Artikel 191). Wir werden nun beweisen, dass $\frac{\mathfrak{A}}{\mathfrak{C}}$ einem der Brüche $\frac{\alpha''}{\gamma''}$, $\frac{\alpha'''}{\gamma'''}$, $\frac{\alpha''''}{\gamma''''}$ usw. gleich ist und $\frac{\mathfrak{B}}{\mathfrak{D}}$ dem unmittelbar folgenden; ist also $\frac{\mathfrak{A}}{\mathfrak{C}}$ $=\frac{\alpha^m}{\gamma^m}$, so ist $\frac{\mathfrak{B}}{\mathfrak{D}}$ $=\frac{\alpha^{m+1}}{\gamma^{m+1}}$. Im vorhergehenden Artikel haben wir gezeigt, dass die Größen $\frac{\alpha'}{\gamma'}$, $\frac{\alpha''}{\gamma''}$, $\frac{\alpha'''}{\gamma'''}$ usw. (die wir kurz mit (1), (2), (3) usw. bezeichnen werden) zusammen mit $L$ die folgende Reihenfolge (I) einhalten: $(1),(3),(5)\mathinner{\ldotp\ldotp\ldotp} L\mathinner{\ldotp\ldotp\ldotp}(6),(4),(2)$. Die erste dieser Größen ist $=0$ (wegen $\alpha'=0$); alle anderen haben dasselbe Vorzeichen wie $L$ oder $a$. Nach Voraussetzung haben aber $\frac{\mathfrak{A}}{\mathfrak{C}}$, $\frac{\mathfrak{B}}{\mathfrak{D}}$ (für die wir $\mathfrak{M},\mathfrak{N}$ schreiben werden) dasselbe Vorzeichen. Offenbar liegen diese Größen also rechts von (1) (oder, wenn man lieber will, auf derselben Seite wie $L$); da $L$ zwischen ihnen liegt, liegt die eine rechts von $L$, die andere links davon. Wir können leicht zeigen, dass $\mathfrak{M}$ nicht rechts von (2) liegen kann. Sonst läge nämlich $\mathfrak{N}$ zwischen (1) und $L$. Daraus würde \textit{erstens} folgen, dass (2) zwischen $\mathfrak{M}$ und $\mathfrak{N}$ liegt und daher der Nenner des Bruchs (2) größer ist als der von $\mathfrak{N}$ (Artikel 190); \textit{zweitens}, dass $\mathfrak{N}$ zwischen (1) und (2) liegt und daher der Nenner von $\mathfrak{N}$ größer ist als der von (2). \textit{Das ist ein Widerspruch.}

Nehmen wir an, $\mathfrak{M}$ sei keinem der Brüche (2), (3), (4) usw. gleich, und sehen wir, was daraus folgt. Liegt der Bruch $\mathfrak{M}$ links von $L$, so muss er offenbar entweder zwischen (1) und (3), zwischen (3) und (5), zwischen (5) und (7) usw. liegen (denn $L$ ist irrational, also sicher von $\mathfrak{M}$ verschieden, und die Brüche (1), (3), (5) usw. können $L$ näher kommen als jede gegebene Größe, die von $L$ verschieden ist). Liegt dagegen $\mathfrak{M}$ rechts von $L$, so liegt es notwendig entweder zwischen (2) und (4), zwischen (4) und (6), zwischen (6) und (8) usw. Nehmen wir also an, $\mathfrak{M}$ liege zwischen
\SourceLeaf{187}
$(m)$ und $(m+2)$. Offenbar liegen dann die Größen $\mathfrak{M},(m),(m+1),(m+2),L$ in der folgenden Reihenfolge:
\begin{center}
(II)\HistoricalNote{*)}{Es spielt hier keine Rolle, ob die Reihenfolge in (II) dieselbe wie in (I) oder die entgegengesetzte ist, \textit{d.h.}, ob $(m)$ auch in (I) links von $L$ oder rechts davon liegt.}:\qquad $(m),\ (\mathfrak{M}),\ (m+2),\ L,\ (m+1)$
\end{center}
Dann muss $\mathfrak{N}=(m+1)$ sein. Denn $\mathfrak{N}$ liegt rechts von $L$. Läge es aber auch rechts von $(m+1)$, so läge $(m+1)$ zwischen $\mathfrak{M}$ und $\mathfrak{N}$, woraus $\gamma^{m+1}>\mathfrak{C}$ folgt; dagegen läge $\mathfrak{M}$ zwischen $(m)$ und $(m+1)$, woraus $\mathfrak{C}>\gamma^{m+1}$ folgt (Artikel 190), \textit{ein Widerspruch}. Läge dagegen $\mathfrak{N}$ links von $(m+1)$, also zwischen $(m+2)$ und $(m+1)$, so wäre $\mathfrak{D}>\gamma^{m+2}$; und weil $(m+2)$ zwischen $\mathfrak{M}$ und $\mathfrak{N}$ liegt, wäre $\gamma^{m+2}>\mathfrak{D}$, \textit{ein Widerspruch}. Also gilt $\mathfrak{N}=(m+1)$ oder $\frac{\mathfrak{B}}{\mathfrak{D}}=\frac{\alpha^{m+1}}{\gamma^{m+1}}=\frac{\beta^m}{\delta^m}$.

Da $\mathfrak{A}\mathfrak{D}-\mathfrak{B}\mathfrak{C}=1$, ist $\mathfrak{B}$ zu $\mathfrak{D}$ teilerfremd, und aus demselben Grund ist $\beta^m$ zu $\delta^m$ teilerfremd. Daher kann die Gleichung $\frac{\mathfrak{B}}{\mathfrak{D}}=\frac{\beta^m}{\delta^m}$ offenbar nur gelten, wenn entweder $\mathfrak{B}=\beta^m$, $\mathfrak{D}=\delta^m$ oder $\mathfrak{B}=-\beta^m$, $\mathfrak{D}=-\delta^m$ ist. Da nun die Form $f$ durch die eigentliche Substitution $\alpha^m,\beta^m,\gamma^m,\delta^m$ in die Form $f^m$ übergeht, die gleich $(\pm a^m,b^m,\mp a^{m+1})$ ist, haben wir die Gleichungen
\begin{align*}
a\alpha^m\alpha^m+2b\alpha^m\gamma^m-a'\gamma^m\gamma^m&=\pm a^m \quad \mathinner{\ldotp\ldotp\ldotp\ldotp\ldotp} \tag*{[5]}\\
a\alpha^m\beta^m+b(\alpha^m\delta^m+\beta^m\gamma^m)-a'\gamma^m\delta^m&=b^m \quad \mathinner{\ldotp\ldotp\ldotp\ldotp} \tag*{[6]}\\
a\beta^m\beta^m+2b\beta^m\delta^m-a'\delta^m\delta^m&=\mp a^{m+1} \quad \mathinner{\ldotp\ldotp\ldotp\ldotp\ldotp} \tag*{[7]}\\
\alpha^m\delta^m-\beta^m\gamma^m&=1 \quad \mathinner{\ldotp\ldotp\ldotp\ldotp\ldotp\ldotp\ldotp\ldotp\ldotp} \tag*{[8]}
\end{align*}
Daraus folgt (aus den Gleichungen 7 und 3) $\mp a^{m+1}=-A'$. Multiplizieren wir ferner Gleichung [2] mit $\alpha^m\delta^m-\beta^m\gamma^m$ und Gleichung [6] mit $\mathfrak{A}\mathfrak{D}-\mathfrak{B}\mathfrak{C}$ und subtrahieren, so bestätigt man durch Ausmultiplizieren leicht, dass
\begin{equation*}
\begin{aligned}
B-b^m={}&(\mathfrak{C}\alpha^m-\mathfrak{A}\gamma^m)\bigl(a\mathfrak{B}\beta^m\\
&\qquad+b(\mathfrak{D}\beta^m+\mathfrak{B}\delta^m)-a'\mathfrak{D}\delta^m\bigr)\\
&+(\mathfrak{B}\delta^m-\mathfrak{D}\beta^m)\bigl(a\mathfrak{A}\alpha^m\\
&\qquad+b(\mathfrak{C}\alpha^m+\mathfrak{A}\gamma^m)-a'\mathfrak{C}\gamma^m\bigr)
\quad \mathinner{\ldotp\ldotp}
\end{aligned}
\tag*{[9]}
\end{equation*}
oder, da entweder $\beta^m=\mathfrak{B}$, $\delta^m=\mathfrak{D}$ oder $\beta^m=-\mathfrak{B}$, $\delta^m=-\mathfrak{D}$,
\[
\begin{aligned}
B-b^m&=\pm(\mathfrak{C}\alpha^m-\mathfrak{A}\gamma^m)(a\mathfrak{B}\mathfrak{B}+2b\mathfrak{B}\mathfrak{D}-a'\mathfrak{D}\mathfrak{D})\\
&=\mp(\mathfrak{C}\alpha^m-\mathfrak{A}\gamma^m)A'
\end{aligned}
\]
Also gilt $B\equiv b^m\ (\mathrm{mod.}\ A')$. Da aber sowohl $B$ als auch $b^m$ zwischen $\surd D$ und $\surd D\mp A'$ liegen, muss $B=b^m$ und daher $\mathfrak{C}\alpha^m-\mathfrak{A}\gamma^m=0$ oder $\frac{\mathfrak{A}}{\mathfrak{C}}=\frac{\alpha^m}{\gamma^m}$, \textit{d.h.}, $\mathfrak{M}=(m)$ sein.

So haben wir aus der Annahme, $\mathfrak{M}$ sei keiner der Größen (2), (3), (4) usw. gleich, hergeleitet, dass es tatsächlich einer von ihnen gleich ist. Nehmen wir dagegen von
\SourceLeaf{188}
Anfang an $\mathfrak{M}=(m)$ an, so gilt offenbar entweder $\mathfrak{A}=\alpha^m$, $\mathfrak{C}=\gamma^m$ oder $-\mathfrak{A}=\alpha^m$, $-\mathfrak{C}=\gamma^m$. In beiden Fällen folgt aus [1] und [5] $A=\pm a^m$ und aus [9] $B-b^m=\pm(\mathfrak{B}\delta^m-\mathfrak{D}\beta^m)A$ oder $B\equiv b^m\ (\mathrm{mod.}\ A)$. Wie oben schließen wir daraus $B=b^m$ und damit $\mathfrak{B}\delta^m=\mathfrak{D}\beta^m$. Weil $\mathfrak{B}$ zu $\mathfrak{D}$ und $\beta^m$ zu $\delta^m$ teilerfremd ist, gilt entweder $\mathfrak{B}=\beta^m$, $\mathfrak{D}=\delta^m$ oder $-\mathfrak{B}=\beta^m$, $-\mathfrak{D}=\delta^m$; also folgt aus [7] $-A'=\mp a^{m+1}$. Daher sind die Formen $F$, $f^m$ identisch. Mithilfe der Gleichung $\mathfrak{A}\mathfrak{D}-\mathfrak{B}\mathfrak{C}=\alpha^m\delta^m-\beta^m\gamma^m$ lässt sich aber leicht beweisen, dass wir $+\mathfrak{B}=\beta^m$, $+\mathfrak{D}=\delta^m$ setzen müssen, wenn $+\mathfrak{A}=\alpha^m$, $+\mathfrak{C}=\gamma^m$ ist; dagegen $-\mathfrak{B}=\beta^m$, ${-\mathfrak{D}=-\delta^m}$, wenn $-\mathfrak{A}=\alpha^m$, $-\mathfrak{C}=\gamma^m$ ist. \textit{Damit ist die Behauptung bewiesen.}

III. Ist das Vorzeichen der Größen $\frac{\mathfrak{A}}{\mathfrak{C}}$ usw. dem von $a$ entgegengesetzt, so ist der Beweis dem vorhergehenden so ähnlich, dass es genügt, seine Hauptpunkte anzugeben. Es liegt $\frac{-\surd D+b}{a'}$ zwischen $\frac{\mathfrak{C}}{\mathfrak{A}}$ und $\frac{\mathfrak{D}}{\mathfrak{B}}$. Der Bruch $\frac{\mathfrak{D}}{\mathfrak{B}}$ ist einem der Brüche
\begin{equation*}
\frac{{}'\delta}{{}'\beta},\ \frac{{}''\delta}{{}''\beta},\ \frac{{}'''\delta}{{}'''\beta}\text{ usw. gleich}
\quad \mathinner{\ldotp\ldotp\ldotp\ldotp\ldotp\ldotp\ldotp} \tag*{(I)}
\end{equation*}
\begin{equation*}
\text{setzen wir diesen }=\frac{{}^m\delta}{{}^m\beta},\qquad \frac{\mathfrak{C}}{\mathfrak{A}}\text{ ist dann }=\frac{{}^m\gamma}{{}^m\alpha}
\quad \mathinner{\ldotp\ldotp\ldotp\ldotp\ldotp\ldotp\ldotp\ldotp\ldotp\ldotp\ldotp} \tag*{(II)}
\end{equation*}
Wir beweisen (I) wie folgt. Nehmen wir an, $\frac{\mathfrak{D}}{\mathfrak{B}}$ sei keinem jener Brüche gleich, so muss es zwischen zwei solchen Brüchen $\frac{{}^m\delta}{{}^m\beta}$ und $\frac{{}^{m+2}\delta}{{}^{m+2}\beta}$ liegen. Daraus folgt wie oben notwendig
\[
\frac{\mathfrak{C}}{\mathfrak{A}}=\frac{{}^{m+1}\delta}{{}^{m+1}\beta}=\frac{{}^m\gamma}{{}^m\alpha}
\]
und entweder $\mathfrak{A}={}^m\alpha$, $\mathfrak{C}={}^m\gamma$ oder $-\mathfrak{A}={}^m\alpha$, $-\mathfrak{C}={}^m\gamma$. Da aber $f$ durch die eigentliche Substitution ${}^m\alpha,{}^m\beta,{}^m\gamma,{}^m\delta$ in die Form
\[
{}^m f=(\pm{}^m a,{}^m b,\mp{}^{m-1}a)
\]
übergeht, ergeben sich drei Gleichungen. Aus diesen folgt zusammen mit den Gleichungen 1, 2, 3, 4 und mit ${}^m\alpha{}^m\delta-{}^m\beta{}^m\gamma=1$ wie oben, dass das erste Glied $A$ der Form $F$ gleich dem ersten Glied von ${}^m f$ ist und das mittlere Glied der ersten Form dem der zweiten modulo $A$ kongruent ist. Da beide Formen reduziert sind und somit das mittlere Glied beider zwischen $\surd D$ und $\surd D\mp A$ liegt, sind diese mittleren Glieder gleich. Daraus folgt $\frac{{}^m\delta}{{}^m\beta}=\frac{\mathfrak{D}}{\mathfrak{B}}$. Wir haben also hier die Richtigkeit der Behauptung (I) aus der Annahme hergeleitet, dass sie falsch sei.

Nehmen wir aber $\frac{{}^m\delta}{{}^m\beta}=\frac{\mathfrak{D}}{\mathfrak{B}}$ an, so beweisen wir auf genau dieselbe Weise und mit denselben Gleichungen,
\SourceLeaf{189}
dass auch ${\frac{{}^m\gamma}{{}^m\beta}}=\frac{\mathfrak{C}}{\mathfrak{A}}$ gilt, was (II) besagte. Mithilfe der Gleichungen $\mathfrak{A}\mathfrak{D}-\mathfrak{B}\mathfrak{C}=1$, ${}^m\alpha{}^m\delta-{}^m\beta{}^m\gamma=1$ folgt daraus, dass entweder
\[
\mathfrak{A}={}^m\alpha,\ \mathfrak{B}={}^m\beta,\ \mathfrak{C}={}^m\gamma,\ \mathfrak{D}={}^m\delta
\]
oder
\[
-\mathfrak{A}={}^m\alpha,\ -\mathfrak{B}={}^m\beta,\ -\mathfrak{C}={}^m\gamma,\ -\mathfrak{D}={}^m\delta
\]
ist und dass die Formen $F$, ${}^m f$ identisch sind. \textit{Damit ist die Behauptung bewiesen.}
